\section{Дифференциальное исчисление функций многих переменных: дополнительные разделы}
\subsection{Неявные функции}
\begin{example}
	Рассмотрим $x^2+y^2=1$. Знаем, что это окружность с центром в т. $(0,0)$ и радиусом 1. Тогда решение: $y=\pm\sqrt{1-x^2}$, заметим, что функцией это не является, так как оно не однозначно.\\
	В этом параграфе далее будет доказано, что, по крайней мере локально, такие уравнения можно решить. \\
	Например, возьмём $(x_0,y_0)\notin{(1,0),(-1,0)}$, которая есть в множестве точек, заданных $x^2+y^2=1$. Тогда существует окрестность $x_0$, в которой каждому $x$ соответствует единственный $y$, то есть существует функция, описывающая пары точек из этой окрестности. Заметим, для вышеуказанных $(1,0)$ и $(-1,0)$ таких окрестностей не существует.
	Положим $F(x,y)=x^2+y^2-1$, тогда $F'_y(1,0)=0$, именно поэтому такой окрестности нет. (сейчас будет доказано)
\end{example}
\begin{note}
	\textit{Достаточно малый} $\eps\lra(\ex\eps_0)(\fa\eps<\eps_0)$
\end{note}
\begin{theorem}\nf{(Теорема о неявной функции, определяемой одним уравнением)}\label{26_th1}\\
	Если $F(\vx,y)$ дифференцируема в некоторой окрестности точки $(\v x_0,y_0)\in\R^{n+1}$, причём $\frac{\vd F}{\vd y}$ непрерывна в этой окрестности, $F(\vx_0,y_0)=0,\ \frac{\vd F}{\vd y}(\vx_0,y_0)\ne0$, то $(\fa\text{ достаточно малого }\eps>0)(\ex\delta>0)$ такая, что в кубе $ K_{\delta,\vx\n}=\prod_{j=1}^n(x_j\n-\delta,x_j\n+\delta)$ единственным образом определена функция $y=\phi(\vx)$ так, что в $K_{\delta,\vx\n}\times(y_0-\eps,y_0+\eps)$ равенства $y=\phi(\vx)$ и $F(\vx, y)=0$ равносильны, при этом $\phi$ дифференцируема в $K_{\delta,\vx\n}$, причём:
	\[\frac{\vd \phi}{\vd x_j}(\vx)=-\frac{\frac{\vd F}{\vd x_j}(\vx,y)}{\frac{\vd F}{\vd y}(\vx,y)}\]
\end{theorem}
\begin{proof}
	Пусть $\Omega$ - шар с центром в $(\vx\n,y_0)$ в котором $F$ дифференцируема, $\frac{\vd F}{\vd y}$ непрерывна и $\frac{\vd F}{\vd y}>0$ для определённости.\\
	Функция $F(\vx\n, y)$ дифференцируема на $(y_0-\eps,y_0+\eps)$, если $\vx_0\times(y_0-\eps,y_0+\eps)\subset\Omega$, причём её производная $>0$.\\ Значит, она строго возрастает, а также $F(\vx\n, y_0+\eps)>0,\ F(\vx\n, y_0-\eps)<0$.\\
	Рассмотрим $F(\vx,y_0+\eps)$ и $F(\vx,y_0-\eps)$, они непрерывны в окрестностях $\vx\n$, следовательно:
	\[(\ex\delta>0)(\fa \vx\in K_{\delta,\vx\n},\ K_{\delta,\vx\n}\times(y_0-\eps,y_0+\eps)\subset\Omega)\ F(\vx, y_0+\eps)>0\text{ и }F(\vx,y_0-\eps)<0\]
	Тогда $F(\vx,y)$ при любых фиксированных $\vx\in K_{\delta,\vx\n}$ является дифференцируемой функцией от $y\in[y_0-\eps,y_0+\eps],\ \frac{\vd F}{\vd y}(\vx, y)>0$. Тогда имеем:
	\[\ex! y\in(y_0-\eps,y_0+\eps):F(\vx,y)=0\lra \phi(\vx)=y\]
	Получили непрерывность в точке $(\vx\n,y_0)$. Теперь докажем дифференцируемость $\phi$ в точке $\vx\n$.\\
	Рассмотрим $\vx\n+\v{\tr x}\in K_{\delta,\vx\n}$, положим $y_0+\tr y:=\phi(\vx\n+\v{\tr x})=y$. \\
	\[\tr F(\vx\n,y_0)=F(\vx\n+\v{\tr x},y_0+\tr y)-F(\vx\n,y_0)=0\]
	При этом, $F$ дифференцируема в $(\vx\n,y_0)$, значит:
	\[\tr F(\vx\n,y_0)=\sum_{j=1}\frac{\vd F}{\vd x_j}(\vx\n,y_0)\tr x_j+\frac{\vd F}{\vd y}(\vx\n,y_0)\tr y+\al(\v{\tr x},\tr y)\]
	При этом:
	\[\al(\v{\tr x},y)=\sigma(|(\v{\tr x},\tr y)|),\ \al_j(\v{\tr x},\tr y):=\frac{\al(\v{\tr x},\tr y)}{|(\v{\tr x}, \tr y)|}\cdot\frac{\tr x_j}{|(\v{\tr x},\tr y)|},\ j=1,\dots,n\]
	\[\beta(\v{\tr x},\tr y):=\frac{\al(\v{\tr x},\tr y)}{|(\v{\tr x},\tr y)|}\cdot\frac{\tr y}{|(\v{\tr x},\tr y)|}\]
	Заметим, $\al(\v{\tr x,\tr y})$ и $\beta(\v{\tr x},\tr y)$ - бесконечно малые при $(\v{\tr x},\tr y)\ra0$, также пусть $\tr y$ - какая-то функция $\psi$ от $\v{\tr x}$. Положим $\t\al(\v{\tr x}):=\al(\v{\tr x}, \psi(\v{\tr x}))$, с $\beta$ аналогично, тогда:
	\begin{multline*}
		\tr F(\vx\n,y_0)=\sum_{j=1}^n\l(\frac{\vd F}{\vd x_j}(\vx\n,y_0)+\al_j(\v{\tr x},\tr y)\r)\tr x_j+\l(\frac{\vd F}{\vd y}(\vx\n,y_0)+\beta(\v{\tr x},\tr y)\r)\tr y=\\
		=\sum_{j=1}^n\l(\frac{\vd F}{\vd x_j}(\vx\n,y_0)+\t\al_j\r)\tr x_j+\l(\frac{\vd F}{\vd y}+\t\beta(\v{\tr x})\r)
	\end{multline*}
	\[
	\tr y=-\sum_{j=1}^n\frac
	{
	\frac{\vd F}{\vd x_j}(\vx\n,y_0)+\t\al_j(\v{\tr x})
	}{
	\frac{\vd F}{\vd y}(\vx\n,y_0)+\t\beta(\v{\tr x})
	}
	\tr x_j
	=\sum_{j=1}^n\l(-\frac
	{
	\frac{\vd F}{\vd x_j}(\vx\n,y_0)
	}{
	\frac{\vd F}{\vd y}(\vx\n,y_0)
	}\r)\tr x_j+\gamma(\v{\tr x})
	\]
	где $\ds\gamma(\v{\tr x})=\sum_{j=1}^n\l(
	\frac
	{
	\frac{\vd F}{\vd x_j}(\vx\n, y_0)
	}{
	\frac{\vd F}{\vd y}(\vx\n,y_0)
	}-\frac{
	\frac{\vd F}{\vd y}(\vx\n,y_0)+\t\al(\v{\tr x})
	}{
	\frac{\vd F}{\vd y}(\vx\n,y_0)+\t\beta(\v{\tr x})
	}\r)\tr x_j$.\\
	Заметим, что $\gamma(\v{\tr x})=\sigma(\v{\tr x})$ при $\v{\tr x}\ra0$.\\
	Значит, $\phi$ дифференцируема в точке $\vx\n$. Аналогичные вещи можно сказать про любые точки, подходящие под условие, а значит требуемое получено.
\end{proof}
\begin{note}
	Теперь докажем аналогичную теорему для нескольких уравнений, точная формулировка будет далее, то есть:
	\begin{equation}\label{26_eq1}
		\begin{cases}
			F_1(x_1,\dots,x_n,y_1,\dots,y_m)=0\\
			\dots\\
			F_m(x_1,\dots,x_n,y_1,\dots,y_m)=0
		\end{cases}
	\end{equation}
	Составим матрицу из производных этих функций по каждому из $y_1,\dots,y_m$:
		\begin{equation}\label{26_m1}
		\begin{pmatrix}
			\frac{\vd F_1}{\vd y_1} & \dots & \frac{\vd F_1}{\vd y_m} \\
			\vdots & \ddots & \vdots \\
			\frac{\vd F_m}{\vd y_1} & \cdots & \frac{\vd F_m}{\vd y_m}
		\end{pmatrix}
	\end{equation}
\end{note}
\begin{definition}
	\textit{Якобианом} $\ds\frac{\mc D(F_1,\dots,F_m)}{\mc D(y_1,\dots,y_m)}$ называется определитель (\ref{26_m1}), где все производные берутся в одной точке.
\end{definition}
\begin{theorem}\nf{(О неявных функциях, определяемых системой уравнений)}\label{26_th2}\\
	Если функции $F_i(\v x,\v y),\ i=1,\dots,m$ дифференцируемы в некоторой окрестности точки $(\vx\n, \v y\n)\in\R^{n+m}$, причём $\frac{\vd F_i}{\vd y_j},\ i,j=1,\dots,m$ непрерывна в этой окрестности, $F_i(\vx\n,\v y\n)=0,\ i=1,\dots,m$, $\frac{\mc D(F_1,\dots,F_m)}{\mc D(y_1,\dots,y_m)}(\vx\n,\v y\n)\ne0$, то $(\fa\text{ достаточно малых }\eps_1>0,\dots,\eps_m>0)(\ex\delta>0)$ такая, что в $K_{\delta,\vx\n}$ существуют и при том единственные функции $\phi_1(\vx),\dots,\phi_m(\vx)$ такие, что $(\fa (\vx,\v y)\in K_{\delta,\vx\n})\times\prod_{i=1}^m(y_i\n-\eps_i,y_i\n+\eps_i)$ система уравнений \normalfont{(\ref{26_eq2})} равносильна \normalfont{(\ref{26_eq1})}. Все функции $\phi_1,\dots,\phi_m$ дифференцируемы в $K_{\delta,\vx\n}$.
	\begin{equation}\label{26_eq2}
		\begin{cases}
			y_1=\phi_1(x_1,\dots,x_n)\\
			\dots\\
			y_m=\phi_m(x_1,\dots,x_n)
		\end{cases}
	\end{equation}
\end{theorem}
\begin{proof}
	Докажем по индукции. Для $m=1$ доказано, см. (\ref{26_th1}). Переход.
	Пока не доказано, на следующей лекции наверно.
\end{proof}